\documentclass[10pt,a4paper]{amsart}
\usepackage{amsmath,amssymb,amsthm,graphicx}
\usepackage{verbatim}
\newenvironment{reference}{\comment}{\endcomment}
%\newenvironment{reference}{}{}
\newenvironment{slogan}{\comment}{\endcomment}
\newenvironment{history}{\comment}{\endcomment}
% For commutative diagrams we use Xy-pic
\usepackage[all]{xy}

% We use 2cell for 2-commutative diagrams.
\xyoption{2cell}
% Theorem environments.
%
\theoremstyle{plain}
\newtheorem{theorem}[subsection]{Theorem}
\newtheorem{proposition}[subsection]{Proposition}
\newtheorem{lemma}[subsection]{Lemma}

\theoremstyle{definition}
\newtheorem{definition}[subsection]{Definition}
\newtheorem{example}[subsection]{Example}
\newtheorem{exercise}[subsection]{Exercise}
\newtheorem{situation}[subsection]{Situation}

\theoremstyle{remark}
\newtheorem{remark}[subsection]{Remark}
\newtheorem{remarks}[subsection]{Remarks}

\numberwithin{equation}{subsection}

% Macros
%
\def\lim{\mathop{\mathrm{lim}}\nolimits}
\def\colim{\mathop{\mathrm{colim}}\nolimits}
\def\Spec{\mathop{\mathrm{Spec}}}
\def\Hom{\mathop{\mathrm{Hom}}\nolimits}
\def\Ext{\mathop{\mathrm{Ext}}\nolimits}
\def\SheafHom{\mathop{\mathcal{H}\!\mathit{om}}\nolimits}
\def\SheafExt{\mathop{\mathcal{E}\!\mathit{xt}}\nolimits}
\def\Sch{\mathit{Sch}}
\def\Mor{\mathop{\mathrm{Mor}}\nolimits}
\def\Ob{\mathop{\mathrm{Ob}}\nolimits}
\def\Sh{\mathop{\mathit{Sh}}\nolimits}
\def\NL{\mathop{N\!L}\nolimits}
\def\CH{\mathop{\mathrm{CH}}\nolimits}
\def\proetale{{pro\text{-}\acute{e}tale}}
\def\etale{{\acute{e}tale}}
\def\QCoh{\mathit{QCoh}}
\def\Ker{\mathop{\mathrm{Ker}}}
\def\Im{\mathop{\mathrm{Im}}}
\def\Coker{\mathop{\mathrm{Coker}}}
\def\Coim{\mathop{\mathrm{Coim}}}

% Boxtimes
%
\DeclareMathSymbol{\boxtimes}{\mathbin}{AMSa}{"02}

%
% Macros for moduli stacks/spaces
%
\def\QCohstack{\mathcal{QC}\!\mathit{oh}}
\def\Cohstack{\mathcal{C}\!\mathit{oh}}
\def\Spacesstack{\mathcal{S}\!\mathit{paces}}
\def\Quotfunctor{\mathrm{Quot}}
\def\Hilbfunctor{\mathrm{Hilb}}
\def\Curvesstack{\mathcal{C}\!\mathit{urves}}
\def\Polarizedstack{\mathcal{P}\!\mathit{olarized}}
\def\Complexesstack{\mathcal{C}\!\mathit{omplexes}}
% \Pic is the operator that assigns to X its picard group, usage \Pic(X)
% \Picardstack_{X/B} denotes the Picard stack of X over B
% \Picardfunctor_{X/B} denotes the Picard functor of X over B
\def\Pic{\mathop{\mathrm{Pic}}\nolimits}
\def\Picardstack{\mathcal{P}\!\mathit{ic}}
\def\Picardfunctor{\mathrm{Pic}}
\def\Deformationcategory{\mathcal{D}\!\mathit{ef}}
\usepackage{cleveref}
\crefname{originalreference}{original reference}{original references}
\Crefname{originalreference}{Original reference}{Original references}
\makeatletter
\newlabel{spaces-morphisms-lemma-permanence-quasi-finite}{{spaces-morphisms-lemma-permanence-quasi-finite}{source}{}{sourcepointer-1}{}}
\newlabel{spaces-morphisms-lemma-permanence-quasi-finite@cref}{{[originalreference][0][]spaces-morphisms-lemma-permanence-quasi-finite}{[1][0][]source}{}{}{}}
\makeatother
\crefname{originalreference}{original reference}{original references}
\Crefname{originalreference}{Original reference}{Original references}
\makeatletter
\newlabel{spaces-morphisms-lemma-permanence-unramified}{{spaces-morphisms-lemma-permanence-unramified}{source}{}{sourcepointer-2}{}}
\newlabel{spaces-morphisms-lemma-permanence-unramified@cref}{{[originalreference][0][]spaces-morphisms-lemma-permanence-unramified}{[1][0][]source}{}{}{}}
\makeatother
\crefname{originalreference}{original reference}{original references}
\Crefname{originalreference}{Original reference}{Original references}
\makeatletter
\newlabel{spaces-morphisms-lemma-unramified-flat-lfp-etale}{{spaces-morphisms-lemma-unramified-flat-lfp-etale}{source}{}{sourcepointer-3}{}}
\newlabel{spaces-morphisms-lemma-unramified-flat-lfp-etale@cref}{{[originalreference][0][]spaces-morphisms-lemma-unramified-flat-lfp-etale}{[1][0][]source}{}{}{}}
\makeatother
\crefname{originalreference}{original reference}{original references}
\Crefname{originalreference}{Original reference}{Original references}
\makeatletter
\newlabel{stacks-properties-lemma-check-property-covering}{{stacks-properties-lemma-check-property-covering}{source}{}{sourcepointer-4}{}}
\newlabel{stacks-properties-lemma-check-property-covering@cref}{{[originalreference][0][]stacks-properties-lemma-check-property-covering}{[1][0][]source}{}{}{}}
\makeatother
\title{Unramifiedness or local quasi-finiteness of a cover forces the same property for the diagonal}
\author{The Stacks Project Authors}
\date{stacks-project:stacks-morphisms.tex:local-snapshot-20261012; GNU FDL 1.2 or later}
\begin{document}
\maketitle
\noindent\textit{Editorial note.} All delivered mathematical text is the author's own, copied from the pinned Stacks Project chapter source sources/stacks-algebraic-stacks-morphisms/pinned-original-source.tex; no statement or proof step was rewritten, reordered or removed.\par\medskip
\noindent\textit{Editorial note.} Preamble provenance: the Stacks Project keeps one single shared preamble for every chapter; this card takes it directly from the staged sources/stacks-algebraic-stacks-morphisms/preamble.tex as three verbatim spans (lines 8-12, 14-18 and 154-217), so no derived file was written and nothing was added or redefined. The wrapper class and the base packages are supplied by the assembler.\par\medskip
\noindent\textit{Editorial note.} Bibliography provenance: neither the counted statement nor its proof contains a citation (this lane rescans the whole span before writing the card), so the card carries no bibliography and no entry was reconstructed.\par\medskip
\noindent\textit{Editorial note.} Selection provenance: this card was opened under the file-based control contract (CONTROL.md), with the mechanical pre-card gate run on its statement and proof interval before the card was written.\par\medskip
\noindent\textit{Allowed standard prerequisites (editorial).}\begin{itemize}
\item The statement that an unramified flat morphism of finite presentation of algebraic spaces is etale. Invoked at the start of the counted proof (source lines 1129-1131) as a cross-chapter pointer. Explains the parenthetical clause of the statement.
\item The criterion for checking a property of a morphism of algebraic stacks along a smooth surjective cover. Invoked inside the counted proof (source lines 1140-1145) as a cross-chapter pointer. Reduces the property of the diagonal to the corresponding property of the comparison morphism.
\item The identification of the pullback of the diagonal with the fibre product of the covering algebraic space with itself. Used inside the counted proof (source lines 1137-1139). Turns the comparison morphism into the canonical map from the fibre product to the square.
\item The permanence statements for unramified and for locally quasi-finite morphisms of algebraic spaces. Invoked at the end of the counted proof (source lines 1146-1155) as cross-chapter pointers. Gives unramifiedness or local quasi-finiteness of the canonical map from the two projections.
\end{itemize}\medskip
\ifdefined\hypertarget\hypertarget{sourcepointer-1}{}\fi
\noindent\textit{Source reference.} Original reference spaces-morphisms-lemma-permanence-quasi-finite: sources/stacks-algebraic-stacks-morphisms/pinned-original-source.tex, source lines 971--971. Navigation pointer only: the counted text refers to the Stacks item spaces-morphisms-lemma-permanence-quasi-finite, which belongs to another chapter of the project and is not part of this card. The pointer carries the target's own label as its number because this chapter's staged snapshot does not contain the other chapter's tag table; its content is not reproduced here.\par\medskip
\ifdefined\hypertarget\hypertarget{sourcepointer-2}{}\fi
\noindent\textit{Source reference.} Original reference spaces-morphisms-lemma-permanence-unramified: sources/stacks-algebraic-stacks-morphisms/pinned-original-source.tex, source lines 968--968. Navigation pointer only: the counted text refers to the Stacks item spaces-morphisms-lemma-permanence-unramified, which belongs to another chapter of the project and is not part of this card. The pointer carries the target's own label as its number because this chapter's staged snapshot does not contain the other chapter's tag table; its content is not reproduced here.\par\medskip
\ifdefined\hypertarget\hypertarget{sourcepointer-3}{}\fi
\noindent\textit{Source reference.} Original reference spaces-morphisms-lemma-unramified-flat-lfp-etale: sources/stacks-algebraic-stacks-morphisms/pinned-original-source.tex, source lines 1131--1131. Navigation pointer only: the counted text refers to the Stacks item spaces-morphisms-lemma-unramified-flat-lfp-etale, which belongs to another chapter of the project and is not part of this card. The pointer carries the target's own label as its number because this chapter's staged snapshot does not contain the other chapter's tag table; its content is not reproduced here.\par\medskip
\ifdefined\hypertarget\hypertarget{sourcepointer-4}{}\fi
\noindent\textit{Source reference.} Original reference stacks-properties-lemma-check-property-covering: sources/stacks-algebraic-stacks-morphisms/pinned-original-source.tex, source lines 479--479. Navigation pointer only: the counted text refers to the Stacks item stacks-properties-lemma-check-property-covering, which belongs to another chapter of the project and is not part of this card. The pointer carries the target's own label as its number because this chapter's staged snapshot does not contain the other chapter's tag table; its content is not reproduced here.\par\medskip
\begin{lemma}
\label{lemma-properties-covering-imply-diagonal}
Let $\mathcal{X}$ be an algebraic stack.
Let $W$ be an algebraic space, and let $f : W \to \mathcal{X}$
be a surjective, flat, locally finitely presented morphism.
\begin{enumerate}
\item If $f$ is unramified (i.e., \'etale, i.e., $\mathcal{X}$
is Deligne-Mumford), then $\mathcal{X}$ is DM.
\item If $f$ is locally quasi-finite, then $\mathcal{X}$ is quasi-DM.
\end{enumerate}
\end{lemma}

\begin{proof}
Note that if $f$ is unramified, then it is \'etale by
Morphisms of Spaces, Lemma
\ref{spaces-morphisms-lemma-unramified-flat-lfp-etale}.
This explains the parenthetical remark in (1).
Assume $f$ is unramified (resp.\ locally quasi-finite). We have to show that
$\Delta_\mathcal{X} : \mathcal{X} \to \mathcal{X} \times \mathcal{X}$
is unramified (resp.\ locally quasi-finite). Note that
$W \times W \to \mathcal{X} \times \mathcal{X}$ is also
surjective, flat, and locally of finite presentation. Hence it suffices to
show that
$$
W \times_{\mathcal{X} \times \mathcal{X}, \Delta_\mathcal{X}} \mathcal{X}
=
W \times_\mathcal{X} W
\longrightarrow
W \times W
$$
is unramified (resp.\ locally quasi-finite), see
Properties of Stacks, Lemma
\ref{stacks-properties-lemma-check-property-covering}.
By assumption the morphism $\text{pr}_i : W \times_\mathcal{X} W \to W$
is unramified (resp.\ locally quasi-finite). Hence
the displayed arrow is unramified (resp.\ locally quasi-finite) by
Morphisms of Spaces, Lemma
\ref{spaces-morphisms-lemma-permanence-unramified}
(resp.\ Morphisms of Spaces, Lemma
\ref{spaces-morphisms-lemma-permanence-quasi-finite}).
\end{proof}

\end{document}
